\chapter{Bits, Bytes, and Buses}

Wick 0.3 grew from a concrete next project: a game in which the player builds
a processor, using the Intel 8080 as the target. A language for that game
needs to make bits visible and register widths explicit. This chapter builds
the small pieces; it does not implement a full CPU or its timing.

\section{Write the bits you mean}

\begin{lstlisting}[language=wick]
let opcode = 0x76
let sign_mask = 0b10000000
let last_address = 0xFFFF
check(opcode == 118, "hexadecimal")
check(sign_mask == 128, "binary")
\end{lstlisting}

Hexadecimal and binary prefixes accept either case. Digits must be valid for
the base and the whole value must fit in 32 unsigned bits. A bare prefix,
an underscore, a fractional prefixed literal, or a value above `0xFFFFFFFF`
is a compile error. Decimal numbers keep their existing rules. All these
values still have type `num`; no integer type was added.

\section{Gates and masks}

The six bit functions are `bit_and`, `bit_or`, `bit_xor`, `bit_not`,
`bit_shl`, and `bit_shr`. The first three combine two values. NOT complements
all 32 bits. Left shift discards bits beyond bit 31; right shift fills with
zero. Operands must be whole numbers in 0 through 4294967295. Shift counts
must be whole numbers in 0 through 31.

\begin{lstlisting}[language=wick]
let wires = 0b10100101
check(bit_and(wires, 0x0F) == 5, "low nibble")
check(bit_xor(wires, 0x80) == 0x25, "toggle bit 7")
check(bit_shr(wires, 7) == 1, "read high bit")
check(bit_shl(1, 31) == 0x80000000, "high bit")
\end{lstlisting}

Type and argument-count mistakes are compile errors. Numeric domain errors,
such as a fractional operand or a shift count of 32, are runtime errors
reported at the Wick call site. The VM validates the value before any C++
integer conversion or shift. NaN and infinity are rejected too.

These are ordinary typed functions. Their addition changes neither operator
precedence nor the meaning of `and`, `or`, and `not`, which still operate on
booleans.

\section{A register has a width}

`u8(n)` wraps a safe integer modulo 256; `u16(n)` wraps it modulo 65536.
Both accept negative values. Inputs must be finite whole numbers between
-9007199254740991 and 9007199254740991 inclusive. Ordinary arithmetic remains
double arithmetic: wrapping happens only where you request it.

\begin{lstlisting}[language=wick]
record Register { value: num, name: str }
let a = Register { value: 0xFF, name: "A" }
let wide = a.value + 1
let carry = wide > 0xFF
a.value = u8(wide)
check(a.value == 0 and carry, "keep the carry")
check(u8(-1) == 255, "byte underflow")
check(u16(0xFFFF + 1) == 0, "PC rollover")
check(u8(bit_not(0x0F)) == 0xF0, "8-bit NOT")
\end{lstlisting}

Keep the wide result until the flags have been calculated. Wrapping first
would erase the very carry the player needs to see. For an eight-bit ADD,
zero is `result == 0`, sign tests mask `0x80`, and auxiliary carry tests
whether the low-nibble sum exceeds `0x0F`. Parity counts the set bits of
the wrapped result and is true for an even count.

\section{Pairs and memory}

A high byte and a low byte make a sixteen-bit address:

\begin{lstlisting}[language=wick]
let h = 0x80
let l = 0x05
let address = bit_or(bit_shl(h, 8), l)
check(address == 0x8005, "register pair")
check(bit_shr(address, 8) == h, "high byte")
check(u8(address) == l, "low byte")
let memory: list<num> = []
for i in 0..65536 { push(memory, 0) }
memory[address] = 0x76
check(memory[0x8005] == 118, "memory write")
\end{lstlisting}

This list models the address space without a new container type. Its entries
are still numbers, so write `u8(value)` at a byte-write boundary. It is not a
packed byte buffer and makes no promise of hardware clock speed. The game
can advance the processor by deliberate simulation steps independently of
render frames.

\section{Make the state readable}

`hex(value, width=1)` returns uppercase digits without a prefix;
`bin(value, width=1)` returns binary digits. Values use the unsigned 32-bit
domain. Width is an integer minimum of 1--8 for hex and 1--32 for binary.
A value is never truncated to fit a label.

\begin{lstlisting}[language=wick]
check(hex(10, 2) == "0A", "byte label")
check(bin(5, 8) == "00000101", "bus label")
check(hex(256, 2) == "100", "overflow stays visible")
\end{lstlisting}

\section{The workbench and the next game}

Run `./build/lantern games/bitlab` to see these operations in a small Wick
workbench. Arrow keys select and edit the input registers; Z cycles ADD,
AND, and XOR. Touch or click an input bit to toggle it. The flags follow
8080 ADD, ANA, and XRA behaviour, including ANA's auxiliary-carry rule:
it reflects bit 3 of either input, rather than the 8085 rule of always
setting that flag. The reference is Intel's \emph{8080/8085 Assembly Language
Programming} manual (1977), chapter 1, ``Auxiliary Carry Flag.''

Bit Lab is an example for the language release. Building the player's
circuit editor, progression, instruction decoder, and complete processor
verification remains game work. No circuit simulator is hidden in the host;
the example's logic is Wick code.

\paragraph{Upgrading.} The ten new built-in names reject same-named user function declarations at compile time. Rename a
colliding user function. No imports, closures, nested records, string indexing, or new
numeric types are introduced in this release.
